\begin{table}[htbp]
\centering
\small
\caption{Illustrative Sentence-Level Communicative Stance}
\label{tab:examples}
\begin{threeparttable}
\begin{tabular}{p{0.13\textwidth}p{0.67\textwidth}c}
\toprule
Class & Illustrative sentence & $\hat p_{\text{class}}$ \\
\midrule
Supportive & ``Inventors are encouraged to consult with the OTT concerning their options regarding the release of intellectual property.'' [Public R1, land-grant, 2012] & 0.973 \\
Neutral & ``Special procedures relevant to these processes are posted on the website of the Office of the Vice President for Research and Innovation.'' [Public R1, 2021] & 0.980 \\
Restrictive & ``As a condition of affiliation with the University, members of the University community are bound by all University policies, including this one.'' [Private R1, 2005] & 0.982 \\
\midrule
\multicolumn{3}{l}{\textit{Same disclosure obligation, different framing}} \\
Supportive & ``The University intends that the obligation for disclosure shall not unduly interfere with prompt publication of research results.'' [Public R1, 2010] & 0.975 \\
Restrictive & ``\ldots\ Each Invention must be disclosed promptly upon its creation, conception or discovery.'' [Public R1, 2013] & 0.978 \\
\bottomrule
\end{tabular}
\begin{tablenotes}[flushleft]\footnotesize
\item \textit{Notes:} $\hat p_{\text{class}}$ is the calibrated model probability assigned to the displayed class. Brackets give the institutional type and source year; institutions are identified in the replication package. The last two rows both state a disclosure obligation. The examples illustrate the annotation construct; they are not an independent accuracy test.
\end{tablenotes}
\end{threeparttable}
\end{table}
